
\section{Findings: Barriers that Discourage Lab and Field \rev{Scientific Practitioners} from Adopting AI in their Workflows}
\label{sec:barriers}

Refer to Figure ~\ref{fig:teaser} on Page 1 to see three barriers that hindered participants from using AI for lab and field work.


\subsection{\rev{Scientific practitioners} do not want to face risk of AI errors since physical experimental setups are too high-stakes}

The most frequently-mentioned barrier to AI usage was the trade-off between potential time saved and the devastating cost of AI-induced errors for physical setups that are hard to setup and maintain. This concern spanned labs both large and small, especially those that worked on projects that required substantial investments of time and money.

For example, P1, a \rev{doctoral student} neuroscience researcher, described the extensive, months-long cycle of his lab's experiments: acquiring and training rodents (spans 6 months), hand-building delicate custom electronic ``drives'' designed to measure neural activity (spans 3 months), surgically implanting these tiny drive devices into rodents, constructing bespoke experimental setups, and running behavioral trials to collect data. In particular, pre-experiment surgeries were especially high-stakes and time-consuming: \textit{``If we are hitting one [brain] region in a surgery we've done before, it can usually take eight hours. The first surgery for [labmate's project] took, I want to say closer to 16 hours, because it was multiple [regions]'' (P1).} In this context, P1 feared relying on AI for informational support because even a single mistake could cost him \textit{``the 10 hours that I was in surgery at that point, but also the months that I had spent training this rat---if the surgery goes wrong and the rat dies.''} P1 summed up the risk in one question: \textit{``How much time does [AI] save versus how much time does it cost me if it's wrong?}'' Figure ~\ref{fig:teaser}A illustrates this example.


%The difficulty of calibrating trust in AI tools is apparent; P1 must actively decide every time at what point the potential saved time outweighs irreversible error.


Participants described similar stakes at larger industrial labs, with an additional element of errors being dangerous and/or amplifying due to AI inaccuracy. P4, P6, and P7's scientific work revolve around designing fuel capsules for experiments at U.S. national labs aiming to achieve nuclear fusion. These experiments rely on rare and costly ``beamtime''~\cite{traweek2009beamtimes} opportunities to  access expensive lab-wide shared laser systems. Collecting sufficient data for even a single experiment can span years, and failed trials cannot easily be repeated (P7). P4 further emphasizes that: \textit{``It's imperative that every data set that we acquire is correct [...] there's a lot more at stake in that measurement. If you're limited to 10 opportunities, if you mess up one, then you can't get that back. Any delays that we incur are going to affect everyone downstream from this as well; in severe cases, that can jeopardize the shot.''}

That is why P4 preferred human judgment and human validation over AI tools: \textit{``Because of the gravity of the data that we're producing, people are much more comfortable trusting another person to do that.''} P6 added that it was hard to trust AI due to its inability to express nuance: \textit{``ChatGPT's writing style makes it sound like it is 100\% correct. Like it'll confidently tell you things that are actually wrong, and that's really dangerous [...] if I'm proposing a hypothesis or giving my opinion, I want to also have the nuance of saying, `Hey, this is what I think, but it might not be entirely true.'''}


\subsection{Physical environments of lab and field sites make it hard to use current AI tools}

Embodied research often takes place in challenging environments that require \rev{participants} to, for example, collect observational data outdoors (P2, P8) or work in highly-controlled settings to protect sensitive experimental conditions (P1, P4, P6), such as in nuclear fusion materials development. Because a substantial portion of their time is spent in these environments, many \rev{participants} could not physically access electronic devices that run AI tools like ChatGPT. 

For instance, during a tour of P7's nuclear fusion lab, she gestured toward a notebook placed outside of lab doors and explained that all external equipment, including phones, laptops, and paper notebooks, were prohibited inside due to contamination risks, which meant no way to access AI assistants inside. Scientists jotted down notes before entering the lab, where full-body clean-room suits are required: hood, hair bonnet, face veil, gloves, coveralls, and boots. P12, who works with cancer cells, must also work in similar clean-room environments (see Figure ~\ref{fig:teaser}B) using laminated paper documentation rather than electronic devices: \textit{``You always have your gloves on with some reagents on it. Maybe some cells are even on it [...] you take [the laminated paper] with you wherever you go, so you can still follow the protocol. At the end of the day, you will spray it, clean it.''} In another instance, P1 noted that he could not have cell phones or laptops in the experiment room because electromagnetic waves, such as Wi-Fi, may interfere with electrode data collection.

In the field AI-enabled devices were avoided not because of contamination risk, but to prevent behavioral interference with the animals being studied such as when brightly-lit devices or voice dictation might distract primates and harm data collection for P8. He remarked that while he did occasionally use his phone to consult ChatGPT, it often was \textit{"messy when you do testing. I'm wearing gloves. I am in a sandy, grimy environment, where I don't want to take out the phone to type. It's definitely a little bit more rugged and reliable to have pen and paper and a clipboard."}

\subsection{AI cannot match the tacit knowledge and contextual judgment of human scientists}

Laboratory and field science involve lots of tacit knowledge consisting of hard-to-capture `unwritten' expertise and embodied decision making~\cite{polanyi2009tacit}, where scientists rely on physical senses such as sight, touch, and even smell to evaluate problems in situ. Many participants were skeptical of AI’s ability to match humans in these skills.

For instance, Figure ~\ref{fig:teaser}C shows how P3 found ChatGPT unable to help troubleshoot mechanical issues with a robot while out at experimental field sites in deserts and volcanos: \textit{``Because it's hard to explain it to the AI in a way that makes sense. Like, I can't really explain how the robot moves; it's difficult to use my words.''}

P10's workplace once debated whether to use AI to automate the placing of fragile organoids beneath a microscope. However, they decided it was more efficient and reliable for a human to do this work, in part because the lab can be a very ``noisy'' environment that makes current AI tools infeasible:

\begin{quote}\textit{``Say you have your pipette and it falls on the ground...it's dirty, but you have something in your hand, so you put your stuff down, you take the pipette up and replace it. It's contaminated. Maybe you should change your gloves. That's so much noise for just one action that happens with something falling on the ground [...] A human brain is actually unmatchable, that biological brain, because we're made exactly for tasks like that [...] We can very easily adjust based on previous knowledge, gut feeling, intuition}.'' (P10)\end{quote}

Moreover, due to the trial-and-error nature of such work, P10 added that some tasks may only be performed once, which makes training a task-specific AI replacement not worthwhile. P10 described this threshold as a \textit{``golden zone''} where it is \textit{``better to let the human do it.''}

A field researcher also shared how his tacit expertise could not be replicated by AI. Over time, P8 has developed the ability to distinguish between chimpanzee subjects by sight: \textit{``For a group of like, 15 chimpanzees, I can tell them apart. Others cannot [...] at some point, you have an intuition.''} Despite attempts to improve computer vision models, AI tools to identify animal subjects \textit{``is still way behind to be good enough for the kind of complex data we have.''}

P12 explained that AI could only assist with piecemeal tasks, such as drafting a protocol for \textit{``specific positions, or like small pieces of experiment planning.''} However, planning an entire experiment end-to-end requires sequencing tasks across multiple physical stations, a process that relies on one's detailed tacit knowledge of their lab's architectural layout accumulated from firsthand experience. As P12 put it, \textit{“First I go into the cell culture... then I do my RNA extraction... then I go to the other room... AI cannot do this [because it] cannot connect the different [lab] stations. It doesn't know the timeline, the context. It cannot plan a whole experiment for me because it doesn't know how an experiment is set up.''}

Lastly, P4 mentioned how evaluating scientific work could be a surprisingly subjective process, which is what separates an expert from a novice and cannot be feasibly replicated by current AI tools. He explained, \textit{``There isn't necessarily a clear right or wrong in terms of data quality [...] is this within the realm of acceptability? Are there negative influences from, let's say, the environment or instrument that would warrant me to have to retake the data set?''} P4 later used the term \textit{``artisan''} to describe the nature of scientific procedures that often evolve over time, which highlights the interpretive, experience-driven aspects of experimental work.